\section{The probabilistic view}
\label{sec:probabilistic}

Every method in these notes was introduced the same way: we chose a function class, we
chose a loss, and we minimised. The choice of loss was always motivated, but always
somewhat by hand. The squared error was ``by far the most common''; the logistic loss
won a beauty contest against three rivals; the $K$-means objective simply asserted that
groups should be tight. The reader is entitled to ask where these expressions actually
come from, and whether there is any principle behind them or only tradition.

There is a principle, and this section is about it. In every case, the objective we
minimised is the \emph{negative log-likelihood} of an explicit probabilistic model of
how the data was generated. Different losses are not competing conventions: they are
different assumptions about the data, written in a different alphabet. This is worth a
section of its own for three reasons. It explains where our losses come from; it tells
us how to build a new one when the standard assumptions do not fit; and it delivers
something the risk-minimisation view never gave us, namely predictions that come with
probabilities attached rather than bare numbers.

\subsection{Maximum likelihood in one page}
\label{subsec:mle}

\paragraph{The principle.}
Suppose we posit a family of probability distributions $p_\theta$, indexed by a
parameter $\theta$, as candidate explanations for how the data was produced. Having
observed data $\mathcal{D}$, the \emph{likelihood} of $\theta$ is the probability (or
density) that the model assigns to what we actually saw,
\[
\theta \ \longmapsto \ p_\theta(\mathcal{D}) .
\]
\emph{Maximum likelihood estimation} (MLE) picks the parameter making the observed data
as unsurprising as possible:
\[
\hat\theta \ \in \ \argmax_{\theta} \ p_\theta(\mathcal{D}) .
\]
Note the logic, which students often reverse: the data is fixed and the parameter
varies. The likelihood is \emph{not} a probability distribution over $\theta$.

\paragraph{Why the logarithm, and why the minus sign.}
If the $n$ samples are independent, the likelihood is a product, and products are
unpleasant: they underflow, and their derivatives are a mess. Since $\log$ is strictly
increasing, maximising $p_\theta(\mathcal{D})$ is the same as maximising $\log
p_\theta(\mathcal{D})$, which turns the product into a sum. Flipping the sign to get a
minimisation, as is our habit, gives the \emph{negative log-likelihood}
\[
\mathrm{NLL}(\theta) \ = \ - \log p_\theta(\mathcal{D})
\ = \ - \sum_{i=1}^n \log p_\theta(\text{sample } i) ,
\]
and dividing by $n$, which does not move the minimiser (the same remark as in
\Cref{sec:regression}), puts it in exactly the shape of an empirical risk:
\[
\boxed{\
\frac{1}{n}\, \mathrm{NLL}(\theta)
\ = \ \frac{1}{n} \sum_{i=1}^n \underbrace{\big( - \log p_\theta(\text{sample } i) \big)}_{\text{a loss } \ell}
\ }
\]
This is the whole dictionary. \emph{A probabilistic model of one sample is a loss
function, and the loss is minus the log of the model's density.} The rest of this
section is that sentence applied three times.

\begin{remark}[Additive constants are free]
Terms in the negative log-likelihood that do not involve $\theta$ can be dropped: they
shift the objective without moving its minimiser. We use this constantly below, and it
is why factors like $(2\pi)^{-d/2}$ never survive into the final loss.
\end{remark}


\subsection{Least squares is maximum likelihood with Gaussian noise}
\label{subsec:mle-ls}

\paragraph{The model.}
Assume, exactly as in \Cref{subsec:bv-ridge}, that the response is a linear function of
the input corrupted by independent Gaussian noise:
\[
y_i = \langle w, x_i \rangle + \varepsilon_i ,
\qquad
\varepsilon_1, \dots, \varepsilon_n \ \text{ i.i.d. } \ \mathcal{N}(0, \sigma^2) .
\]
Here $w$ is the parameter of the \emph{model} we are fitting, and it plays both roles at
once: the model asserts that the truth is $\langle w, \cdot\rangle$ for some $w$, and
maximum likelihood then searches over exactly that $w$. (In \Cref{subsec:bv-ridge}, where
the truth and the fit had to be compared, the two were carefully distinguished as
$w^\star$ and $\hat w$. Here they need not be, since the whole point is to \emph{find}
the parameter, and the estimator we end up with is the $\hat w$ of that subsection.)
Equivalently, the model says that conditionally on the input, the response is Gaussian
around the prediction,
\[
p_w\big( y_i \mid x_i \big)
= \frac{1}{\sqrt{2\pi\sigma^2}} \, \exp\Big( - \frac{\big( y_i - \langle w, x_i\rangle \big)^2}{2\sigma^2} \Big) .
\]

\begin{prop}[Least squares is the MLE]
\label{prop:mle-ls}
Under this model, and for any fixed $\sigma^2 > 0$,
\[
\argmax_{w \in \R^d} \ p_w\big( y \mid X \big)
\ = \ \argmin_{w \in \R^d} \ \frac{1}{2n} \big\Vert X w - y \big\Vert^2 ,
\]
that is, the maximum likelihood estimator of $w$ is exactly the least-squares solution
of \Cref{sec:least-squares}.
\end{prop}

\begin{proof}
By independence the likelihood factorises, so
\[
- \log p_w(y \mid X)
= - \sum_{i=1}^n \log p_w(y_i \mid x_i)
= \frac{1}{2\sigma^2} \sum_{i=1}^n \big( y_i - \langle w, x_i\rangle \big)^2
+ \frac{n}{2} \log\big( 2\pi\sigma^2 \big) .
\]
The second term does not involve $w$, and the first is $\tfrac{n}{\sigma^2}$ times the
least-squares objective $L(w) = \tfrac{1}{2n}\Vert Xw - y\Vert^2$. Multiplying an
objective by a positive constant and adding another does not move its minimisers, so
the two problems have the same solutions.
\end{proof}

Two things are worth extracting from this innocuous computation.

\paragraph{The mysterious $\tfrac12$ was never mysterious.}
The factor $\tfrac12$ we carried through the squared loss since \Cref{sec:regression},
justified there only as ``it cancels the $2$ from the derivative'', is the $\tfrac12$
in the exponent of a Gaussian density. So is the square itself. Once the noise is
declared Gaussian, there is no freedom left: the loss is determined.

\paragraph{The noise level is itself estimable.}
We treated $\sigma^2$ as fixed, but nothing stops us from maximising over it too.
Differentiating the negative log-likelihood above in $\sigma^2$ and cancelling gives
\[
\hat\sigma^2 = \frac{1}{n} \big\Vert X \hat w - y \big\Vert^2 ,
\]
the mean squared residual. This is a genuinely new object: it is an estimate of the
irreducible noise floor $\sigma^2$ of \Cref{sec:bias-variance}, which up to now was an
unknown constant we could only reason about. It is also the ingredient that turns a
prediction into a \emph{predictive distribution}: instead of returning the number
$\langle \hat w, x\rangle$, the model returns $\mathcal{N}(\langle \hat w, x\rangle,
\hat\sigma^2)$, and hence an interval rather than a point. Be warned that $\hat\sigma^2$
is biased downwards, for the reason of \Cref{subsec:cv}: the residuals are computed on
the very data used to fit $\hat w$. The standard correction divides by $n - d$ instead
of $n$.

\begin{remark}[Change the noise, change the loss]
The real payoff is that the dictionary runs both ways. The menu of losses in
\Cref{sec:regression} was presented as a list of options; each entry is in fact a
choice of noise distribution.
\begin{itemize}
    \item Gaussian noise, $p(\varepsilon) \propto e^{-\varepsilon^2/2\sigma^2}$, gives
    the \emph{squared} error.
    \item Laplace noise, $p(\varepsilon) \propto e^{-|\varepsilon|/b}$, gives
    $-\log p = |\varepsilon|/b + \text{const}$, that is, the \emph{absolute} error. Its
    heavier tails make a large residual far less improbable, which is precisely the
    formal content of the vague claim that the absolute error is ``more robust to
    outliers''.
    \item A distribution that is Gaussian in the middle and Laplace in the tails gives
    the \emph{Huber} loss, and the crossover point of the two regimes is the Huber
    parameter.
\end{itemize}
So the question ``which loss should I use?'' becomes the far more answerable question
``how do I expect my data to be wrong?''. This is the main practical benefit of the
probabilistic view.
\end{remark}

\begin{exercise}[$\star$]
Carry out the Laplace computation: write the negative log-likelihood of the model $y_i
= \langle w, x_i\rangle + \varepsilon_i$ with $\varepsilon_i$ i.i.d.\ of density
$\tfrac{1}{2b}e^{-|\varepsilon|/b}$, and check that minimising it is minimising
$\sum_i | y_i - \langle w, x_i\rangle |$.
\end{exercise}

\begin{exercise}[$\star\star$]
Suppose the noise is Gaussian but \emph{heteroscedastic}: $\varepsilon_i \sim
\mathcal{N}(0, \sigma_i^2)$ with known but unequal variances. Show that the MLE now
minimises a \emph{weighted} least-squares objective $\sum_i (y_i - \langle w,
x_i\rangle)^2 / \sigma_i^2$, and interpret the weights.
\end{exercise}


\subsection{Ridge is maximum likelihood with a prior}
\label{subsec:map}

The same computation explains regularisation, which so far we justified only by its
effects. The move is to treat the parameter $w$ itself as random, with a
\emph{prior} distribution encoding what we believe about it before seeing any data. Take
the natural choice, a centred Gaussian with variance $\tau^2$ in every coordinate:
\[
w \ \sim \ \mathcal{N}\big( 0, \tau^2 I_d \big) ,
\qquad\text{i.e.}\qquad
p(w) \propto \exp\Big( - \frac{\Vert w \Vert^2}{2\tau^2} \Big) .
\]
Bayes' rule gives the \emph{posterior} density of $w$ given the data, proportional to
likelihood times prior, and the \emph{maximum a posteriori} (MAP) estimator maximises
it. Taking minus the logarithm turns the product into a sum:
\[
- \log p\big( w \mid X, y \big)
= \underbrace{\frac{1}{2\sigma^2} \big\Vert X w - y \big\Vert^2}_{\text{from the likelihood}}
+ \underbrace{\frac{1}{2\tau^2} \big\Vert w \big\Vert^2}_{\text{from the prior}}
+ \ \text{constant} .
\]

\begin{prop}[Ridge is the MAP estimator]
\label{prop:map-ridge}
The MAP estimator under the Gaussian prior above is the ridge estimator
$w^\star_\lambda$ of \Cref{subsec:ridge}, with
\[
\lambda = \frac{\sigma^2}{n\, \tau^2} .
\]
\end{prop}

\begin{proof}
Divide the displayed objective by $n/\sigma^2$, which does not move its minimiser:
\[
\frac{1}{2n} \big\Vert Xw - y \big\Vert^2 + \frac{\sigma^2}{2 n \tau^2} \Vert w \Vert^2 ,
\]
which is $L_\lambda(w) = L(w) + \tfrac{\lambda}{2}\Vert w\Vert^2$ with $\lambda =
\sigma^2 / (n\tau^2)$.
\end{proof}

The formula for $\lambda$ is the interesting part, because it says what the
regularisation strength \emph{means}: it is the ratio of how noisy we think the data is
to how large we think the weights are. Regularise heavily when the data is noisy
($\sigma^2$ large) or when we are confident the truth is small ($\tau^2$ small); do not
regularise when the prior is flat ($\tau^2 \to \infty$, hence $\lambda \to 0$, recovering
plain least squares). It also explains the $n$ that appears in the ridge closed form
$(X^\top X + n\lambda I)^{-1}X^\top y$: the prior is a fixed amount of information,
while the data accumulates, so its relative weight must decay like $1/n$.

\begin{remark}[Other priors, other penalties]
As with the noise, changing the prior changes the penalty. A Laplace prior $p(w)
\propto \exp(-\Vert w\Vert_1 / b)$ gives the penalty $\Vert w \Vert_1 = \sum_j |w_j|$,
which is the \emph{Lasso}, the standard route to \emph{sparse} solutions in which many
coordinates of $w$ are exactly zero. A flat prior gives no penalty at all. The general
statement is the pleasing symmetry: \emph{the loss encodes the noise, the penalty
encodes the prior}.
\end{remark}

\begin{remark}[MAP is not the whole Bayesian story]
Taking the maximum of the posterior throws away everything except its location. The
fully Bayesian treatment keeps the entire posterior distribution over $w$, and predicts
by averaging over it, which yields error bars on the predictions rather than a single
curve. We do not develop this, but it is worth knowing that ridge is the visible tip of
that construction.
\end{remark}


\subsection{Logistic regression is maximum likelihood with Bernoulli labels}
\label{subsec:mle-logistic}

We now do the same for classification, where the payoff is larger still: the
probabilistic view finally says what the sigmoid \emph{is}, and explains the name
``logistic regression''.

\paragraph{The model.}
For a binary label the only thing to model is the probability of each class. Since a
probability must lie in $[0,1]$ and our model produces a real-valued score $\langle w,
x\rangle$, we need a function squashing $\R$ into $(0,1)$. Take the sigmoid of
\Cref{subsec:logistic}, and declare
\[
\P_w\big( y = +1 \mid x \big) = \sigma\big( \langle w, x \rangle \big) ,
\qquad
\sigma(t) = \frac{1}{1 + e^{-t}} .
\]
(Here $\sigma$ without a square is the sigmoid, not the noise standard deviation of the
previous subsections; see \Cref{subsec:symbols}.) The complementary probability comes
for free from the identity $\sigma(-t) = 1 - \sigma(t)$:
\[
\P_w\big( y = -1 \mid x \big) = 1 - \sigma\big( \langle w, x\rangle \big)
= \sigma\big( - \langle w, x\rangle \big) .
\]
The two cases collapse into a single formula, and this is exactly where the $\pm 1$
convention earns its keep:
\[
\boxed{\ \P_w\big( y \mid x \big) = \sigma\big( y \, \langle w, x \rangle \big) = \sigma(m) \ }
\]
with $m$ the margin of \Cref{subsec:binary}. The margin, introduced there as a
bookkeeping device that packaged correctness into one number, turns out to be the
natural argument of the model's probability.

\begin{prop}[Logistic regression is the MLE]
\label{prop:mle-logistic}
Under this model the negative log-likelihood of the sample is
\[
- \frac{1}{n} \sum_{i=1}^n \log \P_w\big( y_i \mid x_i \big)
= \frac{1}{n} \sum_{i=1}^n \log\Big( 1 + \exp\big( - y_i \langle w, x_i\rangle \big) \Big) ,
\]
which is precisely the logistic risk $L(w)$ of \Cref{subsec:logistic}.
\end{prop}

\begin{proof}
By the boxed identity, $-\log \P_w(y_i \mid x_i) = -\log\sigma(m_i)$, and directly from
the definition of the sigmoid,
\[
- \log \sigma(m) = \log \frac{1}{\sigma(m)} = \log\big( 1 + e^{-m} \big) .
\]
Summing over $i$ and dividing by $n$ gives the claim.
\end{proof}

So the logistic loss was not chosen for its four good properties after all; those
properties are consequences. It is the only loss consistent with modelling the label as
a Bernoulli variable whose parameter is a sigmoid of a linear score.

\paragraph{Why ``regression''.}
The name has puzzled every student who met it, since the method classifies. The answer
is that something \emph{is} being regressed linearly, just not $y$. Invert the model:
\[
\log \frac{\P_w(y = +1 \mid x)}{\P_w(y = -1 \mid x)}
= \log \frac{\sigma(\langle w, x\rangle)}{\sigma(-\langle w, x\rangle)}
= \langle w, x \rangle .
\]
The quantity on the left is the \emph{log-odds}, or \emph{logit}, of the positive class.
Logistic regression is thus an ordinary linear model \emph{for the log-odds}, which is
also why $\langle w, x\rangle$ deserved the name ``score'': it is a log-odds, and the
decision boundary $\langle w, x\rangle = 0$ is the set where the two classes are judged
equally likely.

\paragraph{Multiclass: the softmax weights are probabilities.}
The same reading upgrades \Cref{subsec:multiclass}. Model the label as a categorical
variable whose probabilities are the softmax of the $K$ scores,
\[
\P_W\big( y = k \mid x \big) = p_k(x)
= \frac{\exp(\langle w_k, x\rangle)}{\sum_{k'=1}^K \exp(\langle w_{k'}, x\rangle)} ,
\]
which is legitimate precisely because the $p_k(x)$ are positive and sum to $1$. Then
\[
- \frac1n \sum_{i=1}^n \log \P_W\big( y_i \mid x_i \big)
= \frac1n \sum_{i=1}^n \Big[ - \langle w_{y_i}, x_i\rangle + \log \sum_{k=1}^K \exp\big( \langle w_k, x_i \rangle \big) \Big] ,
\]
which is verbatim the cross-entropy objective of \Cref{subsec:multiclass}. So the
``softmax weights'', which we introduced as a convenient shorthand appearing in a
gradient, are the model's estimated class probabilities, and this is what a call to
\texttt{predict\_proba} returns. It also retrospectively explains the redundancy of the
parametrisation: adding a constant to every score leaves the probabilities unchanged,
because only differences of log-probabilities are identifiable.

\begin{remark}[Why separable data sends $\Vert w\Vert$ to infinity]
\Cref{subsec:separable} showed that on separable data the logistic risk has no
minimiser and $\Vert w_t \Vert \to \infty$, which looked like a technical pathology. The
probabilistic view makes it obvious. Maximising the likelihood means making the
observed labels as probable as possible, and the model can reach probability $1$ only in
the limit $\sigma(m) \to 1$, that is $m \to +\infty$. If a $w$ classifies everything
correctly, the likelihood is still strictly less than $1$ and can always be improved by
scaling $w$ up. The optimiser is not misbehaving: it is being asked for certainty, and
certainty costs infinite log-odds. Adding the $\ell_2$ penalty is, by
\Cref{subsec:map}, exactly imposing a prior that forbids arbitrarily confident models.
\end{remark}

\begin{remark}[A probability is not automatically a good probability]
Reading $\sigma(\langle \hat w, x\rangle)$ as a probability is licensed by the model,
not by reality. If the model is misspecified, the numbers can be systematically
over-confident even when the classification is accurate: of all the points to which a
model assigns probability $0.9$, appreciably fewer than $90\%$ may actually be positive.
A model for which that proportion is right is called \emph{calibrated}, and checking
calibration on held-out data is a separate exercise from measuring accuracy. Modern
neural networks are notoriously badly calibrated in this sense.
\end{remark}

\begin{exercise}[$\star$]
Recover the $\{0,1\}$ form of the objective directly from the model: show that with
$\tilde y \in \{0,1\}$ and $\sigma_i = \sigma(\langle w, x_i\rangle)$, the likelihood of
one sample is $\sigma_i^{\tilde y_i} (1 - \sigma_i)^{1 - \tilde y_i}$, and that its
negative logarithm is the binary cross-entropy of \Cref{subsec:logistic}.
\end{exercise}

\begin{exercise}[$\star\star$]
The gradient of the logistic risk was $\nabla L(w) = -\tfrac1n\sum_i \sigma(-m_i) y_i
x_i$. Show that it can be rewritten as $\tfrac1n \sum_i \big( \hat p_i - \tilde y_i
\big) x_i$, where $\tilde y_i \in \{0,1\}$ is the label and $\hat p_i = \sigma(\langle
w, x_i\rangle)$ the predicted probability of class $1$. Compare with
\Cref{prop:softmax-grad}: in both cases the gradient is driven by ``predicted
probability minus observed label''. Deduce that at a stationary point the predicted
probabilities match the labels on average in every feature direction.
\end{exercise}


\subsection{$K$-means is maximum likelihood with isotropic Gaussians}
\label{subsec:mle-kmeans}

Unsupervised learning fits the same frame, with one twist: there is no $y$, so what the
model must explain is the distribution of the inputs themselves.

\paragraph{The model.}
Suppose each sample was produced by first choosing a cluster and then drawing a point
from a Gaussian centred at that cluster's centroid, with identity covariance. Writing
$z_{i,k} = 1$ when sample $i$ belongs to cluster $k$, as in \Cref{subsec:kmeans},
\[
p\big( x_i \mid \mu, z_i \big) = \mathcal{N}\big( x_i ; \mu_k, I_d \big)
\ \propto \ \exp\Big( - \tfrac{1}{2} \big\Vert x_i - \mu_k \big\Vert^2 \Big)
\qquad \text{when } z_{i,k} = 1 .
\]
The $K$ cases can be folded into a single product, because every exponent but one
carries a factor $z_{i,k} = 0$:
\[
p\big( x_i \mid \mu, z_i \big) \ \propto \ \prod_{k=1}^K
\exp\Big( - \tfrac{1}{2} \big\Vert x_i - \mu_k \big\Vert^2 z_{i,k} \Big) .
\]

\begin{prop}[$K$-means is the MLE]
\label{prop:mle-kmeans}
Under this model, and treating both the centroids $M$ and the assignments $Z$ as
parameters,
\[
- \log \prod_{i=1}^n p\big( x_i \mid \mu, z_i \big)
= \frac{1}{2} \sum_{i=1}^n \sum_{k=1}^K \big\Vert x_i - \mu_k \big\Vert^2 z_{i,k}
\ + \ \text{constant}
\ = \ \frac{1}{2}\, \mathcal{L}(M, Z) + \text{constant} ,
\]
so maximising the likelihood over $(M, Z)$ is minimising the $K$-means objective.
\end{prop}

\begin{proof}
Take the logarithm of the product form above, which turns the product over $i$ and $k$
into a double sum, and collect the normalising constants of the $n$ Gaussian densities
into the additive constant.
\end{proof}

The two limitations of \Cref{subsec:kmeans} are now not defects but assumptions made
visible.
\begin{itemize}
    \item The covariance is $I_d$: \emph{spherical, and identical for every cluster}.
    That is why $K$-means produces round clusters of comparable spread and cuts a long
    thin group in half.
    \item Gaussian tails decay like $e^{-r^2/2}$, so a point far from every centroid is
    assigned an enormous cost. That is why \emph{outliers} drag the centroids so
    effectively. It is the same fragility as the squared error in regression, and for
    literally the same reason.
\end{itemize}
Both point at the same fix: allow a richer model.

\subsection{Gaussian mixtures: summing over the latent variable}
\label{subsec:gmm}

The model above is slightly odd as a probabilistic object, because the assignments
$z_i$ are treated as \emph{parameters} to be optimised, one per sample, growing with the
dataset. The principled alternative is to treat them as unobserved random variables and
\emph{sum them out}. Doing so, and simultaneously relaxing the two assumptions just
identified, gives the \emph{Gaussian mixture model} (GMM).

A sample is drawn by picking a cluster with probability $\pi_k$,
\[
\P(z = k) = \pi_k , \qquad \pi_k \geq 0 , \quad \sum_{k=1}^K \pi_k = 1 ,
\]
and then drawing $x$ from that cluster's own Gaussian, now with its own covariance:
\[
p\big( x \mid z = k \big) = \frac{1}{(2\pi)^{d/2} \det(\Sigma_k)^{1/2}}
\exp\Big( - \tfrac12 (x - \mu_k)^\top \Sigma_k^{-1} (x - \mu_k) \Big) .
\]
Marginalising over the unobserved $z$ gives a density on $\R^d$,
\[
p(x) = \sum_{k=1}^K \pi_k \, p(x \mid z = k) = \sum_{k=1}^K \pi_k \, \mathcal{N}\big( x ; \mu_k, \Sigma_k \big) ,
\]
and fitting the model means minimising the negative log-likelihood
\[
- \sum_{i=1}^n \log p(x_i)
= - \sum_{i=1}^n \log \Big( \sum_{k=1}^K \pi_k \, \mathcal{N}\big( x_i ; \mu_k, \Sigma_k \big) \Big)
\]
over $(\pi_k, \mu_k, \Sigma_k)_{k=1}^K$. Note the shape: a logarithm \emph{of a sum},
which no longer splits. This is what makes the objective non-convex and denies it a
closed form, and it is the continuous counterpart of the combinatorial hardness of
$K$-means.

\paragraph{Soft assignments.}
Since $z$ is now a random variable, one does not assign a point to a cluster; one
computes the posterior probability that it came from each,
\[
\P\big( z = k \mid x_i \big)
= \frac{\pi_k \, \mathcal{N}(x_i ; \mu_k, \Sigma_k)}{\sum_{k'=1}^K \pi_{k'} \, \mathcal{N}(x_i ; \mu_{k'}, \Sigma_{k'})} ,
\]
called the \emph{responsibility} of cluster $k$ for sample $i$. These are numbers in
$(0,1)$ summing to $1$ over $k$: they are the softened version of the row $z_{i,\cdot}$
of the assignment matrix, and structurally they are the softmax weights of
\Cref{subsec:multiclass} once more, a normalised exponential of scores. The model is
fitted by the \emph{expectation-maximisation} algorithm, which alternates computing the
responsibilities with re-estimating $(\pi_k, \mu_k, \Sigma_k)$ as
responsibility-weighted averages. It is Lloyd's algorithm with soft assignments in
place of hard ones, and it comes with the same guarantee and the same weakness: the
likelihood increases at every iteration, and the limit is a local optimum.

\begin{remark}[$K$-means is the zero-temperature limit]
Fix $\Sigma_k = \epsilon I_d$ for every $k$ and $\pi_k = 1/K$, and let $\epsilon \to 0$.
The responsibilities involve ratios of $e^{-\Vert x_i - \mu_k\Vert^2 / 2\epsilon}$,
which tend to $0$ or $1$ according to which centroid is nearest, so the soft
assignments harden and expectation-maximisation degenerates exactly into Lloyd's
algorithm. $K$-means is therefore a Gaussian mixture at zero temperature: the price of
its simplicity is that it must commit to one cluster for every point, however ambiguous
that point may be.
\end{remark}

\begin{exercise}[$\star$]
Two points sit exactly halfway between two centroids. What does $K$-means report for
them, and what does a Gaussian mixture report? Which answer is more informative, and in
what situation would you nonetheless prefer the first?
\end{exercise}


\subsection{Summary, and the limits of the view}
\label{subsec:mle-summary}

Every objective in these notes is now accounted for by a single principle.

\begin{center}
\renewcommand{\arraystretch}{1.35}
\newcolumntype{L}[1]{>{\raggedright\arraybackslash}p{#1}}
\begin{tabular}{@{}L{0.235\textwidth}L{0.35\textwidth}L{0.33\textwidth}@{}}
\toprule
\textbf{Method} & \textbf{Probabilistic model} & \textbf{Resulting objective} \\
\midrule
Least squares & $y \mid x \sim \mathcal{N}(\langle w, x\rangle, \sigma^2)$
  & squared error $\tfrac{1}{2n}\Vert Xw - y\Vert^2$ \\
Absolute / Huber loss & Laplace / mixed noise
  & $\sum_i |y_i - \langle w, x_i\rangle|$, resp.\ Huber \\
Ridge & the above, plus prior $w \sim \mathcal{N}(0, \tau^2 I)$
  & $L(w) + \tfrac{\lambda}{2}\Vert w\Vert^2$, $\lambda = \tfrac{\sigma^2}{n\tau^2}$ \\
Lasso & the above, plus a Laplace prior on $w$
  & $L(w) + \lambda \Vert w \Vert_1$ \\
Logistic regression & $\P(y = +1\mid x) = \sigma(\langle w, x\rangle)$
  & logistic risk $\tfrac1n\sum_i \log(1 + e^{-m_i})$ \\
Softmax regression & $\P(y = k \mid x) = p_k(x)$, categorical
  & cross-entropy \\
$K$-means & $x \sim \mathcal{N}(\mu_k, I_d)$, $k$ a parameter
  & $\Vert X - ZM\Vert_F^2$ \\
Gaussian mixture & $x \sim \sum_k \pi_k \mathcal{N}(\mu_k, \Sigma_k)$, $k$ latent
  & $-\sum_i \log \sum_k \pi_k \mathcal{N}(x_i; \mu_k, \Sigma_k)$ \\
\bottomrule
\end{tabular}
\end{center}

\noindent
Reading the table across, three benefits stand out.
\begin{itemize}
    \item \emph{It explains the losses.} None of the objectives we minimised was
    arbitrary; each encodes an assumption, and knowing the assumption tells you when
    the method will fail. ``Least squares is sensitive to outliers'' and ``Gaussian
    tails make large residuals very improbable'' are the same sentence.
    \item \emph{It generates new methods.} Faced with a problem the standard losses do
    not fit (counts, durations, censored data, heteroscedastic noise), one does not
    guess a loss: one writes down a model of the data and takes minus its logarithm.
    \item \emph{It returns distributions, not just numbers.} A fitted model gives
    $\hat\sigma^2$, predictive intervals, class probabilities and cluster
    responsibilities, none of which the bare minimisation of a loss provides.
\end{itemize}

\begin{remark}[Two honest caveats]
First, \emph{the model is an assumption, not a fact}. The probabilities a fitted model
reports are internal to it, and if the model is wrong they can be confidently wrong; the
calibration warning of \Cref{subsec:mle-logistic} applies throughout. What the
probabilistic view buys is that the assumptions become explicit and therefore
criticisable, which is progress over a loss chosen by habit.

Second, \emph{maximum likelihood is one principle among several}. It says nothing about
generalisation: \Cref{sec:bias-variance} showed that the unbiased estimator, which is
what MLE delivers for least squares, is never the one minimising test error, and that
some shrinkage always helps. The probabilistic view recovers that shrinkage only when a
prior is added by hand, and choosing that prior brings us straight back to
cross-validation. Likelihood explains where our objectives come from; it does not
absolve us from checking how well they predict.
\end{remark}
